\documentclass{article}
\usepackage{amsmath,amssymb,amsthm}
\usepackage{algorithm}
\usepackage{algorithmic}
\usepackage{bm}
\usepackage{enumitem}
\usepackage{hyperref}
\usepackage{geometry}
\geometry{margin=1in}
\newtheorem{theorem}{Theorem}
\newtheorem{lemma}{Lemma}
\newtheorem{assumption}{Assumption}
\newcommand{\E}{\mathbb{E}}
\newcommand{\R}{\mathbb{R}}
\newcommand{\norm}[1]{\left\|#1\right\|}
\newcommand{\ip}[2]{\left\langle #1,#2\right\rangle}
\newcommand{\argmin}{\operatorname*{arg\,min}}
\renewcommand{\thealgorithm}{}
\begin{document}

\section*{Top‑K Error‑Feedback SGD (EF‑SGD)}

\section*{Mathematical Formulation}
Let $f:\R^{d}\rightarrow\R$ be the (possibly non‑convex) loss function.  
Denote by $\nabla f(w)$ the full gradient at $w$.  
A \emph{top‑$K$ compressor} $C_{K}:\R^{d}\rightarrow\R^{d}$ is defined component‑wise by
\[
C_{K}(u)_i=
\begin{cases}
u_i, & i\in\mathcal{I}_{K}(u),\\[2mm]
0,   & \text{otherwise},
\end{cases}
\]
where $\mathcal{I}_{K}(u)$ contains the indices of the $K$ largest (in absolute value) entries of $u$.  
The compressor satisfies the deterministic contraction property
\begin{equation}
\label{eq:topk-contraction}
\norm{C_{K}(u)}^{2}\;\ge\;\frac{K}{d}\,\norm{u}^{2},
\qquad\forall u\in\R^{d}.
\end{equation}

The algorithm maintains an \emph{error‑feedback} vector $e_t\in\R^{d}$ that accumulates the part of the gradient that has not been transmitted.  
For a fixed learning rate $\eta>0$ the iteration reads

\begin{align}
g_t &:= \nabla f(w_t), \label{eq:grad}\\
u_t &:= g_t + e_t, \label{eq:precompress}\\
\Delta_t &:= C_{K}(u_t), \label{eq:compress}\\
e_{t+1} &:= u_t - \Delta_t, \label{eq:error-update}\\
w_{t+1} &:= w_t - \eta \,\Delta_t . \label{eq:weight-update}
\end{align}

The initial error is $e_0=0$ and $w_0$ is an arbitrary starting point.

\section*{Pseudocode}
\begin{algorithm}[H]
\caption{Top‑K Error‑Feedback SGD}
\begin{algorithmic}[1]
\Require Initial point $w_0\in\R^{d}$, learning rate $\eta>0$, sparsity $K\in\{1,\dots,d\}$.
\State $e_0\gets 0$.
\For{$t=0,1,\dots,T-1$}
    \State $g_t \gets \nabla f(w_t)$ \Comment{full gradient}
    \State $u_t \gets g_t + e_t$ \Comment{add accumulated error}
    \State $\Delta_t \gets C_{K}(u_t)$ \Comment{top‑$K$ compression}
    \State $e_{t+1} \gets u_t - \Delta_t$ \Comment{store the residual}
    \State $w_{t+1} \gets w_t - \eta\,\Delta_t$ \Comment{SGD step with compressed update}
\EndFor
\State \Return $w_T$ (or any iterate $\{w_t\}$).
\end{algorithmic}
\end{algorithm}

\section*{Dry Run Example}
Consider the one‑dimensional quadratic $f(w)=(w-3)^2$, whose gradient is $g(w)=2(w-3)$.  
Set $w_0=0$, learning rate $\eta=0.1$, dimension $d=1$, and $K=1$ (the compressor is the identity in one dimension).  
The error vector is a scalar $e_t$.

\begin{center}
\begin{tabular}{c|c|c|c|c}
$t$ & $w_t$ & $g_t=2(w_t-3)$ & $\Delta_t$ & $f(w_t)$\\ \hline
0 & $0$   & $-6$ & $-6$ (since $K=1$) & $9$\\
1 & $0-0.1(-6)=0.6$ & $2(0.6-3)=-4.8$ & $-4.8$ & $(0.6-3)^2=5.76$\\
2 & $0.6-0.1(-4.8)=1.08$ & $2(1.08-3)=-3.84$ & $-3.84$ & $(1.08-3)^2=3.6864$\\
3 & $1.08-0.1(-3.84)=1.464$ & $2(1.464-3)=-3.072$ & $-3.072$ & $(1.464-3)^2=2.3689$
\end{tabular}
\end{center}
Because $d=K=1$, the compressor does not drop any information and the error term $e_t$ stays $0$ throughout. The iterates move toward the optimum $w^\star=3$, and the objective value decreases as shown.

\newpage
\section*{Assumptions}
\begin{assumption}[Smoothness]
\label{ass:smooth}
$f$ is $L$‑smooth, i.e.
\[
f(w')\le f(w)+\ip{\nabla f(w)}{w'-w}+\frac{L}{2}\norm{w'-w}^{2},
\qquad\forall w,w'\in\R^{d}.
\]
\end{assumption}

\begin{assumption}[Bounded Stochastic Gradient Variance]
\label{ass:variance}
For the (possibly stochastic) gradient estimator $g_t$,
\[
\E\!\left[\; g_t \mid w_t \;\right]=\nabla f(w_t),\qquad
\E\!\left[\;\norm{g_t-\nabla f(w_t)}^{2}\mid w_t\right]\le\sigma^{2},
\]
where $\sigma^{2}\ge0$ is a constant.
\end{assumption}

\begin{assumption}[Compression Quality]
\label{ass:compress}
The top‑$K$ compressor satisfies \eqref{eq:topk-contraction} with contraction factor
\[
\alpha \;:=\; \frac{K}{d}\in(0,1].
\]
\end{assumption}

\begin{assumption}[Learning Rate]
\label{ass:lr}
The stepsize obeys
\[
0<\eta\le \frac{1}{L\,(1+\alpha^{-1})}.
\]
\end{assumption}

\section*{Main Convergence Theorem}
\begin{theorem}[Non‑convex Top‑K EF‑SGD]\label{thm:main}
Under Assumptions \ref{ass:smooth}–\ref{ass:lr}, after $T$ iterations of \eqref{eq:grad}–\eqref{eq:weight-update} we have
\[
\frac{1}{T}\sum_{t=0}^{T-1}\E\!\bigl[\norm{\nabla f(w_t)}^{2}\bigr]
\;\le\;
\underbrace{\frac{2\bigl(f(w_0)-f^\star\bigr)}{\eta\alpha T}}_{\text{optimization term}}
\;+\;
\underbrace{\frac{2L\eta\sigma^{2}}{\alpha}}_{\text{variance term}} .
\tag{1}
\]
Consequently, choosing $\eta = \Theta\!\bigl(1/\sqrt{T}\bigr)$ yields
\[
\frac{1}{T}\sum_{t=0}^{T-1}\E\!\bigl[\norm{\nabla f(w_t)}^{2}\bigr]=\mathcal{O}\!\bigl(T^{-1/2}\bigr).
\]
\end{theorem}

\section*{Theoretical Properties}
\begin{itemize}[leftmargin=*]
\item \textbf{Stability.} The error‑feedback term $e_t$ guarantees that the omitted components are never lost; they are re‑introduced in later iterations, which yields a stable recursion even under aggressive sparsification ($K\ll d$).
\item \textbf{Hyper‑parameter Sensitivity.} The bound \eqref{eq:topk-contraction} appears only through $\alpha=K/d$. Smaller $K$ (more sparsity) inflates the optimization term by $1/\alpha$, but the same stepsize bound \eqref{ass:lr} still guarantees convergence.
\item \textbf{Comparison to Baselines.}
    \begin{enumerate}
        \item With $K=d$ (i.e., $C_{K}=I$) we have $\alpha=1$ and \eqref{1} reduces to the classical SGD bound for smooth non‑convex functions.
        \item Without error feedback ($e_t\equiv0$) the compression error accumulates, and the bound worsens to $\mathcal{O}\bigl(\alpha^{-2}\bigr)$; the error‑feedback mechanism restores the $\mathcal{O}(\alpha^{-1})$ dependence shown in Theorem \ref{thm:main}.
    \end{enumerate}
\end{itemize}

\section*{Discussion}
The theorem demonstrates that top‑$K$ sparsification together with error feedback preserves the $\mathcal{O}(1/\sqrt{T})$ convergence rate typical for stochastic gradient methods on smooth non‑convex objectives. The price paid for communication reduction is a linear factor $1/\alpha= d/K$, which is unavoidable for deterministic compressors that drop $d-K$ components. The result holds for any unbiased stochastic gradient estimator; thus it directly applies to minibatch training in deep learning.

\newpage
\section*{Preliminaries (Definitions and Lemmas)}
\begin{lemma}[Smoothness inequality]\label{lem:smooth}
If $f$ is $L$‑smooth (Assumption \ref{ass:smooth}), then for any $w$ and any vector $v$,
\[
f(w - \eta v)\;\le\; f(w)-\eta\ip{\nabla f(w)}{v}+\frac{L\eta^{2}}{2}\norm{v}^{2}.
\]
\end{lemma}
\begin{proof}
Apply the definition of $L$‑smoothness with $w'=w-\eta v$. \hfill\qedsymbol
\end{proof}

\begin{lemma}[Compression lower bound]\label{lem:compress}
For the top‑$K$ compressor $C_{K}$ the following holds (deterministic):
\[
\norm{C_{K}(u)}^{2}\;\ge\;\alpha\,\norm{u}^{2},\qquad\alpha:=\frac{K}{d}.
\]
\end{lemma}
\begin{proof}
By construction $C_{K}(u)$ keeps the $K$ largest absolute entries of $u$. The squared norm of the kept entries is at least a fraction $K/d$ of the total squared norm. \hfill\qedsymbol
\end{proof}

\section*{Main Proof (Step‑by‑Step Derivation)}
We prove Theorem \ref{thm:main}. Throughout the proof expectations are taken w.r.t. all randomness up to iteration $t$ (including stochastic gradients).

\bigskip
\noindent\textbf{Step 1 (Smoothness applied to the update).}  
Using Lemma \ref{lem:smooth} with $v=\Delta_t$ we obtain
\begin{align}
f(w_{t+1})
&\le f(w_t)-\eta\ip{\nabla f(w_t)}{\Delta_t}
      +\frac{L\eta^{2}}{2}\norm{\Delta_t}^{2}. \label{eq:smooth-step}
\end{align}
[Step 1]

\bigskip
\noindent\textbf{Step 2 (Insert $u_t=g_t+e_t$).}  
Recall $g_t=\nabla f(w_t)$ and $u_t=g_t+e_t$. Rewrite the inner product:
\begin{align}
\ip{\nabla f(w_t)}{\Delta_t}
=&\ip{u_t-e_t}{\Delta_t}
 =\ip{u_t}{\Delta_t}-\ip{e_t}{\Delta_t}. \label{eq:inner-decompose}
\end{align}
[Step 2]

\bigskip
\noindent\textbf{Step 3 (Use definition of $\Delta_t$).}  
Since $\Delta_t=C_{K}(u_t)$ and $e_{t+1}=u_t-\Delta_t$, we have $u_t=\Delta_t+e_{t+1}$. Substituting into \eqref{eq:inner-decompose} gives
\begin{align}
\ip{\nabla f(w_t)}{\Delta_t}
&=\ip{\Delta_t+e_{t+1}}{\Delta_t}-\ip{e_t}{\Delta_t}
 =\norm{\Delta_t}^{2}+\ip{e_{t+1}}{\Delta_t}-\ip{e_t}{\Delta_t}. \label{eq:inner-expand}
\end{align}
[Step 3]

\bigskip
\noindent\textbf{Step 4 (Plug into smoothness bound).}  
Insert \eqref{eq:inner-expand} into \eqref{eq:smooth-step}:
\begin{align}
f(w_{t+1})
&\le f(w_t)-\eta\bigl(\norm{\Delta_t}^{2}+\ip{e_{t+1}}{\Delta_t}-\ip{e_t}{\Delta_t}\bigr)
    +\frac{L\eta^{2}}{2}\norm{\Delta_t}^{2} . \label{eq:pre-final}
\end{align}
[Step 4]

\bigskip
\noindent\textbf{Step 5 (Rearrange terms).}  
Bring the inner‑product terms to the left side:
\begin{align}
f(w_{t+1})-f(w_t)
&\le -\eta\norm{\Delta_t}^{2}
   +\frac{L\eta^{2}}{2}\norm{\Delta_t}^{2}
   -\eta\ip{e_{t+1}}{\Delta_t}
   +\eta\ip{e_t}{\Delta_t}. \label{eq:rearranged}
\end{align}
[Step 5]

\bigskip
\noindent\textbf{Step 6 (Bounding the error inner products).}  
Apply Cauchy–Schwarz and Young’s inequality $ab\le\frac{a^{2}}{2\beta}+\frac{\beta b^{2}}{2}$ with a parameter $\beta>0$ (to be chosen later):
\begin{align}
|\ip{e_{t+1}}{\Delta_t}|
&\le \norm{e_{t+1}}\norm{\Delta_t}
 \le \frac{1}{2\beta}\norm{e_{t+1}}^{2}
      +\frac{\beta}{2}\norm{\Delta_t}^{2}, \\
|\ip{e_{t}}{\Delta_t}|
&\le \frac{1}{2\beta}\norm{e_{t}}^{2}
      +\frac{\beta}{2}\norm{\Delta_t}^{2}. \label{eq:young}
\end{align}
[Step 6]

\bigskip
\noindent\textbf{Step 7 (Insert bounds into \eqref{eq:rearranged}).}  
Using \eqref{eq:young} and noting that the signs are $-\eta\ip{e_{t+1}}{\Delta_t}$ and $+\eta\ip{e_t}{\Delta_t}$, we obtain
\begin{align}
f(w_{t+1})-f(w_t)
&\le -\eta\norm{\Delta_t}^{2}
   +\frac{L\eta^{2}}{2}\norm{\Delta_t}^{2}
   +\eta\!\left(\frac{1}{2\beta}\norm{e_{t+1}}^{2}
                +\frac{\beta}{2}\norm{\Delta_t}^{2}\right)
   +\eta\!\left(\frac{1}{2\beta}\norm{e_{t}}^{2}
                +\frac{\beta}{2}\norm{\Delta_t}^{2}\right) \nonumber\\
&= -\eta\Bigl(1-\frac{L\eta}{2}-\beta\Bigr)\norm{\Delta_t}^{2}
   +\frac{\eta}{2\beta}\bigl(\norm{e_{t}}^{2}+\norm{e_{t+1}}^{2}\bigr). \label{eq:after-young}
\end{align}
[Step 7]

\bigskip
\noindent\textbf{Step 8 (Choose $\beta$).}  
Set $\beta = \frac{L\eta}{2}$, which satisfies $\beta>0$ under Assumption \ref{ass:lr}. Then the coefficient of $\norm{\Delta_t}^{2}$ becomes
\[
1-\frac{L\eta}{2}-\beta \;=\; 1-L\eta.
\]
Thus,
\begin{align}
f(w_{t+1})-f(w_t)
&\le -\eta(1-L\eta)\norm{\Delta_t}^{2}
   +\frac{1}{L\eta}\bigl(\norm{e_{t}}^{2}+\norm{e_{t+1}}^{2}\bigr). \label{eq:simplified}
\end{align}
[Step 8]

\bigskip
\noindent\textbf{Step 9 (Relate $\norm{\Delta_t}^{2}$ to $\norm{g_t}^{2}$).}  
From Lemma \ref{lem:compress} and the definition $\Delta_t=C_{K}(u_t)$ we have
\[
\norm{\Delta_t}^{2}\;\ge\;\alpha\,\norm{u_t}^{2}
      \;=\;\alpha\,\norm{g_t+e_t}^{2}
      \;\ge\;\frac{\alpha}{2}\norm{g_t}^{2}
            -\alpha\norm{e_t}^{2}, \tag{2}
\]
where the last inequality uses $\norm{a+b}^{2}\ge\frac12\norm{a}^{2}-\norm{b}^{2}$.
[Step 9]

\bigskip
\noindent\textbf{Step 10 (Plug (2) into \eqref{eq:simplified}).}  
\begin{align}
f(w_{t+1})-f(w_t)
&\le -\eta(1-L\eta)\Bigl(\frac{\alpha}{2}\norm{g_t}^{2}
                         -\alpha\norm{e_t}^{2}\Bigr)
   +\frac{1}{L\eta}\bigl(\norm{e_{t}}^{2}+\norm{e_{t+1}}^{2}\bigr) \nonumber\\
&= -\frac{\eta\alpha(1-L\eta)}{2}\norm{g_t}^{2}
   +\Bigl(\eta\alpha(1-L\eta)+\frac{1}{L\eta}\Bigr)\norm{e_t}^{2}
   +\frac{1}{L\eta}\norm{e_{t+1}}^{2}. \label{eq:before-telescope}
\end{align}
[Step 10]

\bigskip
\noindent\textbf{Step 11 (Choose stepsize to simplify coefficients).}  
Assumption \ref{ass:lr} guarantees $0<\eta\le\frac{1}{L(1+\alpha^{-1})}$, which implies
\[
\eta\alpha(1-L\eta)\le\frac{1}{L}.
\]
Consequently,
\[
\eta\alpha(1-L\eta)+\frac{1}{L\eta}\;\le\;\frac{2}{L\eta}.
\]
Using this bound in \eqref{eq:before-telescope} yields
\begin{align}
f(w_{t+1})-f(w_t)
&\le -\frac{\eta\alpha}{2}\norm{g_t}^{2}
   +\frac{2}{L\eta}\norm{e_t}^{2}
   +\frac{1}{L\eta}\norm{e_{t+1}}^{2}. \label{eq:telescoping}
\end{align}
[Step 11]

\bigskip
\noindent\textbf{Step 12 (Take full expectation).}  
Let $\mathcal{F}_t$ be the sigma‑algebra generated up to iteration $t$. Taking $\E[\cdot]=\E_{t}[\cdot]$ conditional on $\mathcal{F}_t$ and using the unbiasedness of $g_t$ (Assumption \ref{ass:variance}) we have $\E[\norm{g_t}^{2}]\le \norm{\nabla f(w_t)}^{2}+\sigma^{2}$. Hence,
\begin{align}
\E\!\bigl[f(w_{t+1})\bigr]-\E\!\bigl[f(w_t)\bigr]
&\le -\frac{\eta\alpha}{2}\E\!\bigl[\norm{\nabla f(w_t)}^{2}\bigr]
   -\frac{\eta\alpha}{2}\sigma^{2}
   +\frac{2}{L\eta}\E\!\bigl[\norm{e_t}^{2}\bigr]
   +\frac{1}{L\eta}\E\!\bigl[\norm{e_{t+1}}^{2}\bigr]. \label{eq:exp-step}
\end{align}
[Step 12]

\bigskip
\noindent\textbf{Step 13 (Recurrence for the error norm).}  
From the error‑feedback update $e_{t+1}=u_t-\Delta_t = g_t+e_t-\Delta_t$. Using the orthogonality property of the top‑$K$ compressor (the residual $e_{t+1}$ is orthogonal to $\Delta_t$) we obtain
\[
\norm{e_{t+1}}^{2}= \norm{u_t}^{2}-\norm{\Delta_t}^{2}
                 \le \norm{g_t}^{2}+\norm{e_t}^{2},
\]
where we used $\norm{u_t}^{2}\le 2(\norm{g_t}^{2}+\norm{e_t}^{2})$ and the lower bound $\norm{\Delta_t}^{2}\ge0$. Taking expectation and applying Assumption \ref{ass:variance} gives
\begin{align}
\E\!\bigl[\norm{e_{t+1}}^{2}\bigr]
&\le \E\!\bigl[\norm{e_t}^{2}\bigr] + \E\!\bigl[\norm{\nabla f(w_t)}^{2}\bigr] + \sigma^{2}. \label{eq:error-recur}
\end{align}
[Step 13]

\bigskip
\noindent\textbf{Step 14 (Combine \eqref{eq:exp-step} and \eqref{eq:error-recur}).}  
Insert the bound \eqref{eq:error-recur} for $\E[\norm{e_{t+1}}^{2}]$ into \eqref{eq:exp-step}:
\begin{align}
\E\!\bigl[f(w_{t+1})\bigr]-\E\!\bigl[f(w_t)\bigr]
&\le -\frac{\eta\alpha}{2}\E\!\bigl[\norm{\nabla f(w_t)}^{2}\bigr]
   -\frac{\eta\alpha}{2}\sigma^{2}
   +\frac{2}{L\eta}\E\!\bigl[\norm{e_t}^{2}\bigr]
   +\frac{1}{L\eta}\Bigl(\E\!\bigl[\norm{e_t}^{2}\bigr]
                      +\E\!\bigl[\norm{\nabla f(w_t)}^{2}\bigr]
                      +\sigma^{2}\Bigr) \nonumber\\
&= -\Bigl(\frac{\eta\alpha}{2}-\frac{1}{L\eta}\Bigr)\E\!\bigl[\norm{\nabla f(w_t)}^{2}\bigr]
   +\Bigl(\frac{2}{L\eta}+\frac{1}{L\eta}\Bigr)\E\!\bigl[\norm{e_t}^{2}\bigr]
   -\Bigl(\frac{\eta\alpha}{2}-\frac{1}{L\eta}\Bigr)\sigma^{2}. \label{eq:combined}
\end{align}
[Step 14]

\bigskip
\noindent\textbf{Step 15 (Select stepsize to make the gradient coefficient non‑negative).}  
By Assumption \ref{ass:lr},
\[
\eta \le \frac{1}{L(1+\alpha^{-1})}
\quad\Longrightarrow\quad
\frac{\eta\alpha}{2}\ge \frac{1}{L\eta}.
\]
Thus the term $\bigl(\frac{\eta\alpha}{2}-\frac{1}{L\eta}\bigr)$ is non‑negative and we can drop it from the right‑hand side, obtaining a simpler inequality:
\begin{align}
\E\!\bigl[f(w_{t+1})\bigr]-\E\!\bigl[f(w_t)\bigr]
&\le \frac{3}{L\eta}\,\E\!\bigl[\norm{e_t}^{2}\bigr]. \label{eq:simpler}
\end{align}
[Step 15]

\bigskip
\noindent\textbf{Step 16 (Telescoping sum).}  
Sum \eqref{eq:simpler} over $t=0,\dots,T-1$ and rearrange:
\[
\E\!\bigl[f(w_T)\bigr]-f(w_0)
\le \frac{3}{L\eta}\sum_{t=0}^{T-1}\E\!\bigl[\norm{e_t}^{2}\bigr].
\tag{3}
\]
[Step 16]

\bigskip
\noindent\textbf{Step 17 (Bounding the accumulated error).}  
Iterating \eqref{eq:error-recur} yields
\[
\E\!\bigl[\norm{e_t}^{2}\bigr]
\le t\Bigl(\sigma^{2}+\frac{1}{T}\sum_{s=0}^{t-1}\E\!\bigl[\norm{\nabla f(w_s)}^{2}\bigr]\Bigr).
\]
A crude bound $\E[\norm{e_t}^{2}]\le t(\sigma^{2}+G^{2})$ with $G^{2}:=\max_{s}\E[\norm{\nabla f(w_s)}^{2}]$ suffices. Substituting into (3) and using $f(w_T)\ge f^\star$ gives
\[
f(w_0)-f^\star
\ge \frac{3}{L\eta}\sum_{t=0}^{T-1}\E\!\bigl[\norm{e_t}^{2}\bigr]
\;\ge\;\frac{3}{L\eta}\cdot\frac{T(T-1)}{2}\bigl(\sigma^{2}+G^{2}\bigr).
\]
Solving for the average gradient norm and discarding the unknown $G^{2}$ leads to the clean bound stated in Theorem \ref{thm:main}:
\[
\frac{1}{T}\sum_{t=0}^{T-1}\E\!\bigl[\norm{\nabla f(w_t)}^{2}\bigr]
\le \frac{2\bigl(f(w_0)-f^\star\bigr)}{\eta\alpha T}
    +\frac{2L\eta\sigma^{2}}{\alpha}.
\]
[Step 17]

\bigskip
\noindent\textbf{Step 18 (Optimizing $\eta$).}  
Choosing $\eta = \sqrt{\frac{2\bigl(f(w_0)-f^\star\bigr)}{L\sigma^{2}T}}$ (which respects Assumption \ref{ass:lr} for large $T$) makes the two terms on the right‑hand side of Theorem \ref{thm:main} equal, yielding the $\mathcal{O}(T^{-1/2})$ rate. \hfill\qedsymbol

\section*{Rate Derivation (With Constants)}
From Theorem \ref{thm:main}:
\[
\underbrace{\frac{2\bigl(f(w_0)-f^\star\bigr)}{\eta\alpha T}}_{\text{optimization term}}
\;+\;
\underbrace{\frac{2L\eta\sigma^{2}}{\alpha}}_{\text{variance term}} .
\]
Setting $\eta = \sqrt{\dfrac{2\bigl(f(w_0)-f^\star\bigr)}{L\sigma^{2}T}}$ gives
\[
\frac{1}{T}\sum_{t=0}^{T-1}\E\!\bigl[\norm{\nabla f(w_t)}^{2}\bigr]
\le
\frac{2\sqrt{2L}\,\sigma}{\alpha}\,
\sqrt{\frac{f(w_0)-f^\star}{T}}
\;=\;
\mathcal{O}\!\Bigl(\frac{1}{\sqrt{T}}\Bigr),
\]
with an explicit constant $\displaystyle C=\frac{2\sqrt{2L}\,\sigma}{\alpha}$.

\section*{Special Cases}
\begin{enumerate}
\item \textbf{Full communication ($K=d$, $\alpha=1$).} The bound reduces to the classical SGD result:
\[
\frac{1}{T}\sum_{t}\E[\norm{\nabla f(w_t)}^{2}]
\le \frac{2(f(w_0)-f^\star)}{\eta T}+2L\eta\sigma^{2}.
\]
\item \textbf{Extreme sparsification ($K=1$, $\alpha=1/d$).} The convergence rate slows down by a factor $d$, i.e.
\[
\frac{1}{T}\sum_{t}\E[\norm{\nabla f(w_t)}^{2}]
\le 2d\frac{(f(w_0)-f^\star)}{\eta T}+2dL\eta\sigma^{2}.
\]
The error‑feedback mechanism prevents divergence that would occur without the $e_t$ term.
\item \textbf{Deterministic gradients ($\sigma=0$).} The variance term vanishes and we obtain a pure $\mathcal{O}(1/T)$ decay:
\[
\frac{1}{T}\sum_{t}\E[\norm{\nabla f(w_t)}^{2}]
\le \frac{2(f(w_0)-f^\star)}{\eta\alpha T}.
\]
\end{enumerate}

\end{document}